\section{Limitations and Discussion}
\label{sec:discussion}

\paragraph{What is not checked.} The report covers connectivity, interference along single-joint
sweeps from home, mass and a short class test. It does not check multi-joint self-collision over
the whole configuration space, structural strength, manufacturability beyond 3D-printing infill,
or cable routing. Masses depend on density and infill models fitted per robot.

\paragraph{Fidelity is measured against simulation models.} Our references are public MJCF/URDF
models, themselves approximations; agreement with them is not agreement with the hardware. Masses
for parts we do not model (fasteners, wiring) explain most of the remaining gaps (e.g.\ Microduck's
legs are 10\% light).

\paragraph{Examples and library grew together.} The library was extended while the designs were
written, each robot adding what it needed. Q1 therefore shows the
language can express these robots, not that an agent finds them unaided; that is what Q2 measures.

\paragraph{Construction guarantees have edges.} Running the independent checker
(\cref{app:checker}) on our own designs found three faults in the G1 that \ms's checks had passed:
cross-piece joints whose parts overlapped by 5--7.5\,cm$^3$ inside an over-generous pin allowance,
hands drawn 29\,mm from the wrists they belong to (shapes added to a limb's last link were never
joined to it), and three structures that were invalid solids. Each was a gap in a library routine,
not in a design, so each was fixed once in the library (cross pieces are now clevis joints, and
every limb link must be a single solid), after which all eight designs pass both checkers. The
guarantee that nothing floats holds for what the connection API places; derived geometry needs
its own checks, and an independent checker is how we found where they were missing.

\paragraph{Generality.} Designs are limited to rigid bodies with revolute joints, pins and rotors.
Prismatic joints, tendons and soft parts are future work.
